%% wip-stamp.tex
%% ----------------------------------------------------------------------
%% WIP build stamp: prints "[date] [time] [commit hash]" top-left on a
%% page, so every WIP copy of the thesis is traceable to the exact build
%% and commit it came from.
%%
%% Mechanism: a Lua block (package `luacode`, already loaded in
%% preamble.tex) queries the build environment directly.
%%   - date/time: `os.date`, always available, no shell-escape needed.
%%   - commit hash: prefers the CI-provided `GITHUB_SHA` env var (also no
%%     shell-escape needed); falls back to `git rev-parse` via `io.popen`
%%     for local builds, which DOES require shell-escape.
%%
%% Local build: compile with `lualatex -shell-escape` (already set in
%% .latexmkrc). GitHub Actions: `latexmk_shell_escape: true` is set on
%% the xu-cheng/latex-action step in .github/workflows/blank.yml, and
%% GITHUB_SHA is exported into the build container regardless.
%%
%% To disable for good (e.g. the final defense build), remove the
%% %% wip-stamp.tex
%% ----------------------------------------------------------------------
%% WIP build stamp: prints "[date] [time] [commit hash]" top-left on a
%% page, so every WIP copy of the thesis is traceable to the exact build
%% and commit it came from.
%%
%% Mechanism: a Lua block (package `luacode`, already loaded in
%% preamble.tex) queries the build environment directly.
%%   - date/time: `os.date`, always available, no shell-escape needed.
%%   - commit hash: prefers the CI-provided `GITHUB_SHA` env var (also no
%%     shell-escape needed); falls back to `git rev-parse` via `io.popen`
%%     for local builds, which DOES require shell-escape.
%%
%% Local build: compile with `lualatex -shell-escape` (already set in
%% .latexmkrc). GitHub Actions: `latexmk_shell_escape: true` is set on
%% the xu-cheng/latex-action step in .github/workflows/blank.yml, and
%% GITHUB_SHA is exported into the build container regardless.
%%
%% To disable for good (e.g. the final defense build), remove the
%% %% wip-stamp.tex
%% ----------------------------------------------------------------------
%% WIP build stamp: prints "[date] [time] [commit hash]" top-left on a
%% page, so every WIP copy of the thesis is traceable to the exact build
%% and commit it came from.
%%
%% Mechanism: a Lua block (package `luacode`, already loaded in
%% preamble.tex) queries the build environment directly.
%%   - date/time: `os.date`, always available, no shell-escape needed.
%%   - commit hash: prefers the CI-provided `GITHUB_SHA` env var (also no
%%     shell-escape needed); falls back to `git rev-parse` via `io.popen`
%%     for local builds, which DOES require shell-escape.
%%
%% Local build: compile with `lualatex -shell-escape` (already set in
%% .latexmkrc). GitHub Actions: `latexmk_shell_escape: true` is set on
%% the xu-cheng/latex-action step in .github/workflows/blank.yml, and
%% GITHUB_SHA is exported into the build container regardless.
%%
%% To disable for good (e.g. the final defense build), remove the
%% %% wip-stamp.tex
%% ----------------------------------------------------------------------
%% WIP build stamp: prints "[date] [time] [commit hash]" top-left on a
%% page, so every WIP copy of the thesis is traceable to the exact build
%% and commit it came from.
%%
%% Mechanism: a Lua block (package `luacode`, already loaded in
%% preamble.tex) queries the build environment directly.
%%   - date/time: `os.date`, always available, no shell-escape needed.
%%   - commit hash: prefers the CI-provided `GITHUB_SHA` env var (also no
%%     shell-escape needed); falls back to `git rev-parse` via `io.popen`
%%     for local builds, which DOES require shell-escape.
%%
%% Local build: compile with `lualatex -shell-escape` (already set in
%% .latexmkrc). GitHub Actions: `latexmk_shell_escape: true` is set on
%% the xu-cheng/latex-action step in .github/workflows/blank.yml, and
%% GITHUB_SHA is exported into the build container regardless.
%%
%% To disable for good (e.g. the final defense build), remove the
%% \input{wip-stamp.tex} + \WIPStampPage lines in main.tex, or compile
%% with: latexmk -usepretex="\def\NoWIPStamp{}" main.tex
%% ----------------------------------------------------------------------

\begin{luacode*}
  local function trim(s)
    return (s or ""):gsub("^%s+", ""):gsub("%s+$", "")
  end

  local function short_hash(s, n)
    s = trim(s)
    if s == "" then return nil end
    return s:sub(1, n or 7)
  end

  local function git_commit_hash()
    -- CI: GitHub Actions always exports GITHUB_SHA into the build
    -- container; reading an env var needs no shell-escape.
    local sha = short_hash(os.getenv("GITHUB_SHA"), 7)
    if sha then return sha end
    -- Local build: ask git directly (requires -shell-escape).
    local ok, handle = pcall(io.popen, "git rev-parse --short=7 HEAD 2>/dev/null")
    if ok and handle then
      local result = short_hash(handle:read("*a"))
      handle:close()
      if result then return result end
    end
    return "unknown"
  end

  tex.sprint(string.format("\\gdef\\WIPCommitHash{%s}", git_commit_hash()))
  tex.sprint(string.format("\\gdef\\WIPBuildDate{%s}", os.date("%Y-%m-%d")))
  tex.sprint(string.format("\\gdef\\WIPBuildTime{%s}", os.date("%H:%M")))
\end{luacode*}

% Inline stamp text. Braces form a real TeX group so the font/color
% declarations don't leak into the rest of the document.
\newcommand{\WIPStamp}{%
  {\sffamily\footnotesize\color{gray}%
  \WIPBuildDate\ \ \WIPBuildTime\ \ \texttt{\WIPCommitHash}}%
}

% Full page: page style empty, stamp pinned top-left, rest left blank.
% \ifdefined\NoWIPStamp lets a build opt out (see header comment) while
% still leaving a genuinely blank page, matching the pre-stamp layout.
\newcommand{\WIPStampPage}{%
  \ifdefined\NoWIPStamp
    \null\thispagestyle{empty}%
  \else
    \thispagestyle{empty}%
    \noindent\WIPStamp\par
    \vfill\null
  \fi
}
 + \WIPStampPage lines in main.tex, or compile
%% with: latexmk -usepretex="\def\NoWIPStamp{}" main.tex
%% ----------------------------------------------------------------------

\begin{luacode*}
  local function trim(s)
    return (s or ""):gsub("^%s+", ""):gsub("%s+$", "")
  end

  local function short_hash(s, n)
    s = trim(s)
    if s == "" then return nil end
    return s:sub(1, n or 7)
  end

  local function git_commit_hash()
    -- CI: GitHub Actions always exports GITHUB_SHA into the build
    -- container; reading an env var needs no shell-escape.
    local sha = short_hash(os.getenv("GITHUB_SHA"), 7)
    if sha then return sha end
    -- Local build: ask git directly (requires -shell-escape).
    local ok, handle = pcall(io.popen, "git rev-parse --short=7 HEAD 2>/dev/null")
    if ok and handle then
      local result = short_hash(handle:read("*a"))
      handle:close()
      if result then return result end
    end
    return "unknown"
  end

  tex.sprint(string.format("\\gdef\\WIPCommitHash{%s}", git_commit_hash()))
  tex.sprint(string.format("\\gdef\\WIPBuildDate{%s}", os.date("%Y-%m-%d")))
  tex.sprint(string.format("\\gdef\\WIPBuildTime{%s}", os.date("%H:%M")))
\end{luacode*}

% Inline stamp text. Braces form a real TeX group so the font/color
% declarations don't leak into the rest of the document.
\newcommand{\WIPStamp}{%
  {\sffamily\footnotesize\color{gray}%
  \WIPBuildDate\ \ \WIPBuildTime\ \ \texttt{\WIPCommitHash}}%
}

% Full page: page style empty, stamp pinned top-left, rest left blank.
% \ifdefined\NoWIPStamp lets a build opt out (see header comment) while
% still leaving a genuinely blank page, matching the pre-stamp layout.
\newcommand{\WIPStampPage}{%
  \ifdefined\NoWIPStamp
    \null\thispagestyle{empty}%
  \else
    \thispagestyle{empty}%
    \noindent\WIPStamp\par
    \vfill\null
  \fi
}
 + \WIPStampPage lines in main.tex, or compile
%% with: latexmk -usepretex="\def\NoWIPStamp{}" main.tex
%% ----------------------------------------------------------------------

\begin{luacode*}
  local function trim(s)
    return (s or ""):gsub("^%s+", ""):gsub("%s+$", "")
  end

  local function short_hash(s, n)
    s = trim(s)
    if s == "" then return nil end
    return s:sub(1, n or 7)
  end

  local function git_commit_hash()
    -- CI: GitHub Actions always exports GITHUB_SHA into the build
    -- container; reading an env var needs no shell-escape.
    local sha = short_hash(os.getenv("GITHUB_SHA"), 7)
    if sha then return sha end
    -- Local build: ask git directly (requires -shell-escape).
    local ok, handle = pcall(io.popen, "git rev-parse --short=7 HEAD 2>/dev/null")
    if ok and handle then
      local result = short_hash(handle:read("*a"))
      handle:close()
      if result then return result end
    end
    return "unknown"
  end

  tex.sprint(string.format("\\gdef\\WIPCommitHash{%s}", git_commit_hash()))
  tex.sprint(string.format("\\gdef\\WIPBuildDate{%s}", os.date("%Y-%m-%d")))
  tex.sprint(string.format("\\gdef\\WIPBuildTime{%s}", os.date("%H:%M")))
\end{luacode*}

% Inline stamp text. Braces form a real TeX group so the font/color
% declarations don't leak into the rest of the document.
\newcommand{\WIPStamp}{%
  {\sffamily\footnotesize\color{gray}%
  \WIPBuildDate\ \ \WIPBuildTime\ \ \texttt{\WIPCommitHash}}%
}

% Full page: page style empty, stamp pinned top-left, rest left blank.
% \ifdefined\NoWIPStamp lets a build opt out (see header comment) while
% still leaving a genuinely blank page, matching the pre-stamp layout.
\newcommand{\WIPStampPage}{%
  \ifdefined\NoWIPStamp
    \null\thispagestyle{empty}%
  \else
    \thispagestyle{empty}%
    \noindent\WIPStamp\par
    \vfill\null
  \fi
}
 + \WIPStampPage lines in main.tex, or compile
%% with: latexmk -usepretex="\def\NoWIPStamp{}" main.tex
%% ----------------------------------------------------------------------

\begin{luacode*}
  local function trim(s)
    return (s or ""):gsub("^%s+", ""):gsub("%s+$", "")
  end

  local function short_hash(s, n)
    s = trim(s)
    if s == "" then return nil end
    return s:sub(1, n or 7)
  end

  local function git_commit_hash()
    -- CI: GitHub Actions always exports GITHUB_SHA into the build
    -- container; reading an env var needs no shell-escape.
    local sha = short_hash(os.getenv("GITHUB_SHA"), 7)
    if sha then return sha end
    -- Local build: ask git directly (requires -shell-escape).
    local ok, handle = pcall(io.popen, "git rev-parse --short=7 HEAD 2>/dev/null")
    if ok and handle then
      local result = short_hash(handle:read("*a"))
      handle:close()
      if result then return result end
    end
    return "unknown"
  end

  tex.sprint(string.format("\\gdef\\WIPCommitHash{%s}", git_commit_hash()))
  tex.sprint(string.format("\\gdef\\WIPBuildDate{%s}", os.date("%Y-%m-%d")))
  tex.sprint(string.format("\\gdef\\WIPBuildTime{%s}", os.date("%H:%M")))
\end{luacode*}

% Inline stamp text. Braces form a real TeX group so the font/color
% declarations don't leak into the rest of the document.
\newcommand{\WIPStamp}{%
  {\sffamily\footnotesize\color{gray}%
  \WIPBuildDate\ \ \WIPBuildTime\ \ \texttt{\WIPCommitHash}}%
}

% Full page: page style empty, stamp pinned top-left, rest left blank.
% \ifdefined\NoWIPStamp lets a build opt out (see header comment) while
% still leaving a genuinely blank page, matching the pre-stamp layout.
\newcommand{\WIPStampPage}{%
  \ifdefined\NoWIPStamp
    \null\thispagestyle{empty}%
  \else
    \thispagestyle{empty}%
    \noindent\WIPStamp\par
    \vfill\null
  \fi
}
